\section{Introduction}
\label{sec:introduction}

Modern vision models are trained on a fixed source distribution and deployed on a target distribution that inevitably drifts. Even strong ImageNet-pretrained transformers suffer large accuracy drops when the test data is corrupted by noise, blur, weather, or digital artifacts~\citep{hendrycks2019robustness,dosovitskiy2021vit}. \emph{Test-Time Adaptation (TTA)} addresses this gap by adapting the model using only unlabeled data observed at test time, without access to the original source data or any labels~\citep{wang2021tent,sun2020ttt,schneider2020adabn}. Because the target stream is unlabeled and non-stationary, TTA methods must extract a usable learning signal from the model's own predictions while simultaneously avoiding catastrophic failure modes such as confirmation bias, class collapse, and over-confidence.

A recent line of work, culminating in COME (Conservative Opinion-Minimization Entropy)~\citep{come2024come}, advances TTA by grounding the adaptation objective in \emph{evidential deep learning}. COME places a Dirichlet distribution over the class-probability simplex, derives a per-sample \emph{uncertainty mass} $u_i = K / S_i$ from the Dirichlet strength $S_i$, and combines the standard softmax entropy with an \emph{opinion entropy} defined over the belief masses plus the uncertainty mass. The resulting conservative entropy loss down-weights the contribution of uncertain samples, providing a degree of robustness that pure entropy minimization~\citep{wang2021tent} lacks. Empirically, COME improves over Tent~\citep{wang2021tent}, EATA~\citep{niu2022eata}, and SAR~\citep{niu2023sar} on standard corruption benchmarks.

Despite its successes, COME has three structural limitations that motivate our work.

\paragraph{Limitation 1: Uniform treatment of samples.}
COME applies the \emph{same} conservative entropy objective to every test sample, differing only in the magnitude of the per-sample gradient through the (implicit) entropy weighting. There is no mechanism to route genuinely confident samples to a stronger supervisory signal. When the model is already confident and correct on a sample, conservative entropy provides a vanishing gradient, wasting a sample that could otherwise reinforce the decision boundary through a sharper pseudo-label cross-entropy objective.

\paragraph{Limitation 2: No pseudo-label leverage.}
Pseudo-labeling has driven large gains in semi-supervised learning~\citep{sohn2020fixmatch,xie2020selftraining,lee2013pseudo} and has been shown to provide stronger gradients than entropy minimization when the model is confident~\citep{prabhu2023lada,chen2023minimize}. COME does not leverage pseudo-labels at all, leaving a known source of signal on the table. A natural question is whether one can \emph{combine} the safety of conservative entropy (for uncertain samples) with the strength of pseudo-label cross-entropy (for confident samples) within a single TTA framework.

\paragraph{Limitation 3: No calibration guarantee.}
COME uses evidential uncertainty to reweight the loss but never enforces any agreement between the evidential uncertainty $u_i$ and the softmax confidence $c_i = \max_k p_{ik}$ that actually drives predictions. The two quantities can drift arbitrarily far apart during adaptation: the model can become highly confident in its softmax output while the Dirichlet uncertainty remains large, or vice versa. This decoupling undermines both the reliability of the uncertainty estimate (used for error detection) and the trustworthiness of the confidence estimate (used for routing and selective prediction).

\paragraph{Our insight.}
The two quantities that COME computes---softmax confidence $c_i$ and Dirichlet uncertainty $u_i$---are complementary measures of the model's epistemic state. \emph{Confidence} captures the sharpness of the predictive distribution, while \emph{uncertainty} captures the total evidence mass. A sample that is \emph{both} confident ($c_i \to 1$) \emph{and} certain ($u_i \to 0$) is a reliable candidate for pseudo-label supervision; a sample that is uncertain or low-confidence should be handled conservatively. The key realization is that these two measures should be \emph{balanced} to route samples, and that this balance should be \emph{calibrated} so the two views of the model's epistemic state remain consistent.

\paragraph{Our method: UCDPA.}
We propose \textbf{Uncertainty-Calibrated Dual-Path Adaptation (UCDPA)}, which operationalizes this insight through four components:
\begin{enumerate}[leftmargin=*,topsep=2pt,itemsep=2pt]
  \item A novel \textbf{UCB (Uncertainty-Confidence Balance) score} $b_i = \sigma((c_i-\tau_c)/T_c)\,\sigma((1-u_i-\tau_u)/T_u)$ that quantifies, per sample, the degree to which the model is \emph{both} confident and certain. We prove (Theorem~\ref{thm:ucb}) that $b_i$ is monotonically increasing in $c_i$, monotonically decreasing in $u_i$, and achieves its maximum only when $c_i \to 1$ and $u_i \to 0$.
  \item A \textbf{scale-normalized dual-path adaptation loss} that softly routes each sample between conservative opinion-entropy minimization (Path 1, for uncertain/low-confidence samples) and pseudo-label cross-entropy (Path 2, for confident/certain samples), with per-batch scale normalization ensuring both paths contribute comparable gradient magnitudes, weighted by a class-balanced inverse-frequency term from an exponential moving average (EMA) class tracker.
  \item A \textbf{calibration regularizer} $L_\mathrm{cal} = |\E[c] - \E[1-u]|$ that explicitly binds the batch-mean softmax confidence and the batch-mean evidential uncertainty. We show (Proposition~\ref{prop:cal}) that minimizing this regularizer bounds the mean discrepancy between the two uncertainty views, with residual per-sample discrepancy controlled by their joint variance.
  \item A \textbf{feature prototype alignment} term that maintains an EMA prototype bank and pulls features toward their pseudo-label class prototype, improving representation coherence under shift.
\end{enumerate}

\paragraph{Contributions.} Our contributions are:
\begin{enumerate}[leftmargin=*,topsep=2pt,itemsep=2pt]
  \item \textbf{A novel UCB score} that quantifies the balance between Dirichlet uncertainty and softmax confidence for per-sample routing, with a formal monotonicity theorem (Theorem~\ref{thm:ucb}) establishing that the routing probability is monotonically increasing in confidence and monotonically decreasing in uncertainty.
  \item \textbf{A dual-path adaptation framework} that leverages both conservative entropy minimization (for uncertain samples) and pseudo-label cross-entropy (for confident samples), combined with class-balanced reweighting that prevents class collapse and a calibration regularizer that maintains batch-mean consistency between the two uncertainty views (Proposition~\ref{prop:cal}).
  \item \textbf{Comprehensive experiments} on ImageNet-C with ViT-B/16 showing a \ucdpadelta{} percentage-point accuracy improvement over the COME baseline (\ucdpaacc{} vs.\ \comeacc{}), with statistical significance ($p < \pvalue{}$), improved calibration (ECE \ucdpaece{} vs.\ \comeece{}; MCE \ucdpamce{} vs.\ \comemce{}), and improved uncertainty quantification (Predictive Entropy, Mutual Information) and error detection (AUROC \ucdpaauroc{} vs.\ \comeauroc{}). We compare with six recent baselines including CoTTA, MEMO, POEM, and FOA, showing UCDPA outperforms all by $+3.51$--$11.42$~pp. We also validate on CIFAR-100-C with a CIFAR-ResNet-18 backbone (cross-backbone, cross-dataset), and provide sensitivity sweeps (4 routing hyperparameters), batch-size experiments ($B \in \{1, 8, 32, 64\}$, including $B=1$ TTA), a scale-normalization ablation, and a theoretical analysis of routing stability, error amplification, and convergence.
\end{enumerate}

The remainder of the paper is organized as follows. Section~\ref{sec:related} surveys related work. Section~\ref{sec:method} develops the UCDPA method, including the UCB score, dual-path loss, calibration regularizer, prototype alignment, and the monotonicity theorem and consistency proposition. Section~\ref{sec:experiments} reports experiments on ImageNet-C. Section~\ref{sec:discussion} discusses why the method works, its limitations, and threats to validity. Section~\ref{sec:conclusion} concludes.
